\scanpage{44}
\mspage{21}
\[
(53)\qquad\widetilde C'=\widetilde A'\setminus\widetilde B',\qquad C'=A'\setminus B'=\widetilde C'/\sigma_{\widetilde A'},
\]
on écrira
\[
(54)\qquad G=G'/\widetilde C'\quad\text{ou}\quad G=G'/C'
\]
et on dira que $G$ \uline{est déduit de} $G'$ en contractant les arêtes $a'\in C'$ (ou encore, les arcs $\vec a'\in\widetilde C'$) en des pts. Interprétant les morphismes généralisés de graphes, en termes des réalisations \deleted{géométriques}\inserted{topologiques} de ceux-ci, à l'aide d'applications continues qui contractent juste certaines arêtes en des pts, on voit que cette terminologie correspond bien à une intuition topologique correcte. D'ailleurs la somme amalgamée (52) existe quel que soit le sous-graphe $G'_\bullet$ de $G'$ ayant $\unsure{=^{\mathrm s}|\overline m^{\mathrm s}}$ sommets que $G'$, i.e. quel que soit l'ens d'\deleted{flèches}\inserted{arêtes} $C'\subset A'$- de $G'_{\deleted{\unsure{\psi|\alpha}}\bullet}$. [Il s'ensuit que tout morphisme généralisé de graphes $G'\longrightarrow G$ se factorise de façon unique en un
\mspage{43}
